\chapter{Isometric Immersions}
\label{chap:dc-isometric-immersions}
An isometric immersion $f:M^n\to\overline M^{n+m}$ induces the metric
$\langle v,w\rangle=\langle df(v),df(w)\rangle$ on $M$.
\section{The second fundamental form}

Locally identify an immersed patch with its image in $\overline M$. Write $T_p\overline M=T_pM\oplus(T_pM)^\perp$, and decompose fields as $X=X^T+X^N$.

\begin{definition}[the identified setup: tangent distribution and orthogonal splitting]
  \label{def:dc-ch6-2-patch}
  \lean{Riemannian.DCImmersedPatch.normalSpace, Riemannian.DCImmersedPatch.mem_normalSpace_iff, Riemannian.DCImmersedPatch.normalProj, Riemannian.DCImmersedPatch.IsTangentField, Riemannian.DCImmersedPatch.IsNormalField, Riemannian.DCImmersedPatch.isTangentField_bracketField}
  \leanok
  Let $(\overline{M},\langle\,,\rangle)$ be a Riemannian manifold. An
  \emph{immersed patch} of dimension $n$ in $\overline{M}$ — the data furnished
  by an isometric immersion $f:M^n\to\overline{M}$ after the identification of
  $U\subset M$ with $f(U)\subset\overline{M}$ — consists of a family of
  subspaces $T_pM\subseteq T_p\overline{M}$, $p\in\overline{M}$, of constant
  dimension $n$, such that:

  (i) the orthogonal splitting $T_p\overline{M}=T_pM\oplus(T_pM)^\perp$ is
  \emph{differentiable}: every differentiable vector field $X$ on
  $\overline{M}$ admits a differentiable field $X^T$ with $X^T(p)\in T_pM$ and
  $\langle v,\,X(p)-X^T(p)\rangle=0$ for every $v\in T_pM$ and every $p$;

  (ii) fields tangent to $M$ are closed under the Lie bracket: if
  $X(p),Y(p)\in T_pM$ for all $p$, then $[X,Y](p)\in T_pM$ for all $p$ (on $M$,
  $[\overline{X},\overline{Y}]=[X,Y]$).

  The \emph{normal space} at $p$ is
  $(T_pM)^\perp=\{v\in T_p\overline{M};\ \langle w,v\rangle=0\ \text{for all}\
  w\in T_pM\}$, and the \emph{normal part} of a field is $X^N=X-X^T$, so that
  $X=X^T+X^N$. A field is \emph{tangent} when $X(p)\in T_pM$ for every $p$ and
  \emph{normal} when $X(p)\in(T_pM)^\perp$ for every $p$; under the
  identification, the tangent fields are the local fields $\mathcal{X}(U)$ of
  $M$ and the normal fields are $\mathcal{X}(U)^\perp$.
\end{definition}

\begin{lemma}[uniqueness and tensoriality of the tangential projection]
  \label{lem:dc-ch6-2-proj}
  \lean{Riemannian.DCImmersedPatch.eq_zero_of_mem_tang_of_mem_normalSpace, Riemannian.DCImmersedPatch.tangentProj_eq_at, Riemannian.DCImmersedPatch.tangentProj_congr_apply, Riemannian.DCImmersedPatch.tangentProj_add, Riemannian.DCImmersedPatch.tangentProj_smul, Riemannian.DCImmersedPatch.tangentProj_apply_of_mem, Riemannian.DCImmersedPatch.tangentProj_apply_of_mem_normalSpace, Riemannian.DCImmersedPatch.normalProj_add, Riemannian.DCImmersedPatch.normalProj_smul}
  \uses{def:dc-ch6-2-patch}
  \leanok
  The splitting $X(p)=X^T(p)+X^N(p)$ at a point is the unique decomposition of
  $X(p)$ into a tangent and a normal vector: if $u\in T_pM$ and
  $\langle v,\,X(p)-u\rangle=0$ for all $v\in T_pM$, then $u=X^T(p)$. In
  particular $X^T(p)$ depends only on the value $X(p)$; the projections
  $X\mapsto X^T$ and $X\mapsto X^N$ are additive and
  $\mathcal{D}(\overline{M})$-homogeneous; $X^T=X$ and $X^N=0$ for tangent
  fields; and $X^T=0$, $X^N=X$ for normal fields.
\end{lemma}
\begin{proof}
  \leanok
  \sketch
  If $d\in T_pM$ is also normal at $p$ then $\langle d,d\rangle=0$, so $d=0$ by
  positive-definiteness: the splitting is direct. Let $u$ be as in the
  statement. Then $X^T(p)-u\in T_pM$, and for every tangent $v$,
  $\langle v,\,X^T(p)-u\rangle
  =\langle v,\,X(p)-u\rangle-\langle v,\,X(p)-X^T(p)\rangle=0$, so $X^T(p)-u$
  is also normal, hence zero. Every remaining assertion is an instance of this
  uniqueness: $X^T(p)+Y^T(p)$ is tangent and
  $(X+Y)(p)-(X^T(p)+Y^T(p))=(X(p)-X^T(p))+(Y(p)-Y^T(p))$ is orthogonal to
  $T_pM$, so $(X+Y)^T=X^T+Y^T$; likewise $f(p)X^T(p)$ is tangent with
  $f(p)X(p)-f(p)X^T(p)=f(p)(X(p)-X^T(p))$ orthogonal to $T_pM$, so
  $(fX)^T=fX^T$; if $X(p)\in T_pM$ then $u=X(p)$ satisfies the
  characterization, and if $X(p)$ is normal then $u=0$ does. The assertions
  for $X^N=X-X^T$ follow by subtraction.
\end{proof}

\begin{definition}[the induced connection]
  \label{def:dc-ch6-2-induced-connection}
  \lean{Riemannian.DCImmersedPatch.inducedCov}
  \uses{def:dc-ch6-2-patch, thm:dc-ch2-3-6}
  \leanok
  Let $\overline{\nabla}$ be the Riemannian connection of $\overline{M}$. The
  \emph{induced connection} of the patch is
  $\nabla_X Y=(\overline{\nabla}_X Y)^T$.
\end{definition}

\begin{proposition}[the induced connection is the Riemannian connection of the induced metric]
  \label{prop:dc-ch6-2-induced-connection}
  \lean{Riemannian.DCImmersedPatch.isTangentField_inducedCov, Riemannian.DCImmersedPatch.inducedCov_add_left, Riemannian.DCImmersedPatch.inducedCov_smul_left, Riemannian.DCImmersedPatch.inducedCov_add_right, Riemannian.DCImmersedPatch.inducedCov_smul_right, Riemannian.DCImmersedPatch.inducedCov_sub_swap, Riemannian.DCImmersedPatch.inducedCov_compat}
  \uses{def:dc-ch6-2-induced-connection, lem:dc-ch6-2-proj}
  \leanok
  The induced connection is tangent-valued and, on tangent fields, satisfies
  the axioms of an affine connection: it is additive in both arguments,
  $\mathcal{D}$-homogeneous in the direction ($\nabla_{fX}Y=f\nabla_XY$), and
  satisfies the Leibniz rule $\nabla_X(fY)=f\nabla_XY+X(f)\,Y$ for $Y$ tangent.
  Moreover it is symmetric, $\nabla_XY-\nabla_YX=[X,Y]$ for tangent $X,Y$, and
  compatible with the metric on tangent fields,
  $X\langle Y,Z\rangle=\langle\nabla_XY,Z\rangle+\langle Y,\nabla_XZ\rangle$.
  By the uniqueness of the Riemannian connection (\cref{thm:dc-ch2-3-6}) it is
  the Riemannian connection of the induced metric.
\end{proposition}
\begin{proof}
  \leanok
  \sketch
  Additivity and $\mathcal{D}$-homogeneity in the direction slot follow from
  the corresponding laws of $\overline{\nabla}$ and of the projection
  (\cref{lem:dc-ch6-2-proj}). For the Leibniz rule,
  $\overline{\nabla}_X(fY)=f\overline{\nabla}_XY+X(f)\,Y$; taking tangential
  parts, the term $X(f)\,Y$ is already tangent, giving
  $\nabla_X(fY)=f\nabla_XY+X(f)\,Y$. For symmetry, the difference
  $\overline{\nabla}_XY-\overline{\nabla}_YX=[X,Y]$ is tangent for tangent
  $X,Y$ (\cref{def:dc-ch6-2-patch}(ii)), so applying the projection gives
  $\nabla_XY-\nabla_YX=[X,Y]$. For compatibility, write
  $\overline{\nabla}_XY=(\overline{\nabla}_XY)^T+(\overline{\nabla}_XY)^N$ in
  $X\langle Y,Z\rangle
  =\langle\overline{\nabla}_XY,Z\rangle+\langle Y,\overline{\nabla}_XZ\rangle$;
  the normal parts pair to zero against the tangent fields $Z$ and $Y$.
\end{proof}

\begin{definition}[the second fundamental form field]
  \label{def:dc-ch6-2-B}
  \lean{Riemannian.DCImmersedPatch.secondFundForm, Riemannian.DCImmersedPatch.cov_eq_inducedCov_add_secondFundForm, Riemannian.DCImmersedPatch.secondFundForm_mem}
  \uses{def:dc-ch6-2-induced-connection}
  \leanok
  The \emph{second fundamental form} of the patch is
  $B(X,Y)=\overline{\nabla}_XY-\nabla_XY=(\overline{\nabla}_XY)^N$, a
  normal-valued form on fields, giving the \emph{Gauss decomposition}
  $\overline{\nabla}_XY=\nabla_XY+B(X,Y)$.
\end{definition}

\begin{proposition}[the second fundamental form mapping $B$]
  \label{prop:dc-ch6-2-1}
  \lean{Riemannian.DCImmersedPatch.secondFundForm_add_left, Riemannian.DCImmersedPatch.secondFundForm_add_right, Riemannian.DCImmersedPatch.secondFundForm_sub_left, Riemannian.DCImmersedPatch.secondFundForm_smul_left, Riemannian.DCImmersedPatch.secondFundForm_smul_right, Riemannian.DCImmersedPatch.secondFundForm_symm, Riemannian.DCImmersedPatch.isNormalField_secondFundForm}
  \uses{def:dc-ch6-2-B, lem:dc-ch6-2-proj, prop:dc-ch6-2-induced-connection}
  \leanok
  \dcref{ch6:2.1}
  If $X,Y\in\mathcal{X}(U)$, the mapping
  $B:\mathcal{X}(U)\times\mathcal{X}(U)\to\mathcal{X}(U)^\perp$ given by

  $$
  B(X,Y)=\overline{\nabla}_{\overline{X}}\overline{Y}-\nabla_X Y
  $$

  is bilinear and symmetric.
\end{proposition}
\begin{proof}
  \leanok
  \sketch
  From linearity of a connection, $B$ is additive in $X$ and $Y$ and
  $B(fX,Y)=fB(X,Y)$, $f\in\mathcal{D}(U)$. To show $B(X,fY)=fB(X,Y)$, denote the
  extension of $f$ to $\overline{U}$ by $\overline{f}$; then

  $$
  B(X,fY)=\overline{\nabla}_{\overline{X}}(\overline{f}\,\overline{Y})-\nabla_X(fY)=\overline{f}\,\overline{\nabla}_{\overline{X}}\overline{Y}-f\nabla_X Y+\overline{X}(\overline{f})\overline{Y}-X(f)Y.
  $$

  Since $f=\overline{f}$, $\overline{X}(\overline{f})=X(f)$, and $Y=\overline{Y}$ on
  $M$, the last two terms cancel, giving $B(X,fY)=fB(X,Y)$. To show $B$ is
  symmetric, by symmetry of the Riemannian connection,

  $$
  B(X,Y)=\overline{\nabla}_{\overline{X}}\overline{Y}-\nabla_X Y=\overline{\nabla}_{\overline{Y}}\overline{X}+[\overline{X},\overline{Y}]-\nabla_Y X-[X,Y].
  $$

  Since $[\overline{X},\overline{Y}]=[X,Y]$ on $M$, we conclude $B(X,Y)=B(Y,X)$.
\end{proof}

\begin{lemma}[pointwise locality of $B$ and of the covariant derivative]
  \label{lem:dc-ch6-2-locality}
  \lean{Riemannian.AffineConnection.cov_congr_apply_left, Riemannian.exists_decomposition_of_apply_eq_zero, Riemannian.DCImmersedPatch.exists_tangent_decomposition_of_apply_eq_zero, Riemannian.DCImmersedPatch.secondFundForm_congr_apply, Riemannian.DCImmersedPatch.exists_isTangentField_eq, Riemannian.DCImmersedPatch.exists_isNormalField_eq}
  \uses{prop:dc-ch6-2-1, def:dc-ch6-2-patch}
  \leanok
  Because $B$ is bilinear over the differentiable functions, the value
  $B(X,Y)(p)$ depends only on the values $X(p)$ and $Y(p)$: if $X(p)=X'(p)$ and
  the tangent fields $Y,Y'$ satisfy $Y(p)=Y'(p)$, then
  $B(X,Y)(p)=B(X',Y')(p)$. Likewise $\overline{\nabla}_XZ(p)$ depends on the
  direction only through $X(p)$. Moreover every tangent vector $v\in T_pM$
  (resp.\ normal vector $\eta\in(T_pM)^\perp$) is the value at $p$ of a tangent
  (resp.\ normal) field, so the pointwise objects are
  well-defined.
\end{lemma}
\begin{proof}
  \leanok
  \sketch
  A field $\sigma$ with $\sigma(p)=0$ can be written, near $p$, as a finite sum
  $\sum_i f_i W_i$ with differentiable functions $f_i$ vanishing at $p$: expand
  $\sigma$ in a local frame of the tangent bundle of $\overline{M}$ around $p$
  and cut the coefficients and the frame off by a bump function equal to $1$
  near $p$; the discrepancy is a field vanishing on a neighbourhood of $p$.
  For a mapping $\Phi$ that is additive and $\mathcal{D}$-homogeneous in
  $\sigma$, each summand contributes $f_i(p)\,\Phi(W_i)(p)=0$, and a field
  vanishing near $p$ contributes $0$ because it equals $(1-\chi)\tau$ for a
  bump $\chi$ with $\chi(p)=1$ supported in the vanishing neighbourhood.
  Applying this to $\Phi=\overline{\nabla}_{(\cdot)}Z$ and, in the second slot,
  to $\Phi=B(X,\cdot)$ — where for tangent $\sigma$ the frame decomposition
  may be projected tangentially, so the $W_i$ can be taken tangent — gives the
  locality claims. For the extensions: any vector extends to a global
  differentiable field, and composing with the tangential (resp.\ normal)
  projection fixes the value at $p$ when it is tangent (resp.\ normal).
\end{proof}

By \cref{lem:dc-ch6-2-locality}, the value $B(X,Y)(p)$ depends only on the
values $X(p)$ and $Y(p)$. Now let $p\in M$ and $\eta\in(T_pM)^\perp$. The
mapping $H_\eta:T_pM\times T_pM\to\mathbb{R}$ given by

$$
H_\eta(x,y)=\langle B(x,y),\eta\rangle,\quad x,y\in T_pM,
$$

is, by Proposition 2.1, a symmetric bilinear form.

\begin{definition}[second fundamental form $\mathit{II}_\eta$]
  \label{def:dc-ch6-2-2}
  \lean{Riemannian.DCImmersedPatch.secondFundFormAt, Riemannian.DCImmersedPatch.secondFundScalarAt, Riemannian.DCImmersedPatch.secondFundQuadAt, Riemannian.DCImmersedPatch.shapeOperatorAt, Riemannian.DCImmersedPatch.inner_shapeOperatorAt, Riemannian.DCImmersedPatch.inner_shapeOperatorAt_symm, Riemannian.DCImmersedPatch.secondFundScalarAt_symm, Riemannian.DCImmersedPatch.secondFundFormAt_symm}
  \uses{lem:dc-ch6-2-locality, prop:dc-ch6-2-1, def:dc-ch6-3-normal-connection, lem:dc-ch6-3-weingarten}
  \leanok
  \dcref{ch6:2.2}
  The quadratic form $\mathit{II}_\eta$ defined on $T_pM$ by
  $$
  \mathit{II}_\eta(x)=H_\eta(x,x)
  $$
  is called \emph{the second fundamental form} of $f$ at $p$ along the normal vector
  $\eta$. Sometimes the expression \emph{second fundamental form} also designates the
  mapping $B$. The bilinear mapping $H_\eta$ is associated to a linear self-adjoint
  operator $S_\eta:T_pM\to T_pM$ by
  $$
  \langle S_\eta(x),y\rangle=H_\eta(x,y)=\langle B(x,y),\eta\rangle.
  $$
\end{definition}

\begin{proposition}[the operator $S_\eta$ via the covariant derivative]
  \label{prop:dc-ch6-2-3}
  \lean{Riemannian.DCImmersedPatch.shapeOperatorAt_eq_shapeOperator}
  \uses{def:dc-ch6-2-2, def:dc-ch6-3-normal-connection, lem:dc-ch6-3-weingarten, lem:dc-ch6-2-locality}
  \leanok
  \dcref{ch6:2.3}
  Let $p\in M$, $x\in T_pM$ and $\eta\in(T_pM)^\perp$. Let $N$ be a local extension
  of $\eta$ normal to $M$. Then

  $$
  S_\eta(x)=-(\overline{\nabla}_x N)^T.
  $$
\end{proposition}
\begin{proof}
  \leanok
  \sketch
  Let $y\in T_pM$ and let $X,Y$ be local extensions of $x,y$, tangent to
  $M$. Then $\langle N,Y\rangle=0$, and therefore
  
  $$
  \langle S_\eta(x),y\rangle=\langle B(X,Y)(p),N\rangle=\langle\overline{\nabla}_X Y-\nabla_X Y,N\rangle(p)=\langle\overline{\nabla}_X Y,N\rangle(p)=-\langle Y,\overline{\nabla}_X N\rangle(p)=\langle-\overline{\nabla}_x N,y\rangle,
  $$
  
  for all $y\in T_pM$.
\end{proof}

\begin{example}[hypersurface; principal curvatures; Gauss spherical mapping]
  \label{ex:dc-ch6-2-4}
  \lean{Riemannian.gaussMap, Riemannian.gaussMap_norm, Riemannian.normalSpace_eq_span_of_unit_normal, Riemannian.mem_tang_of_inner_unit_normal_eq_zero, Riemannian.fderiv_gaussMap_mem_tang, Riemannian.fderiv_gaussMap_eq_neg_shapeOperatorAt, Riemannian.fderiv_gaussMap_mem_sphereTang}
  \uses{def:dc-ch6-2-2, prop:dc-ch6-2-3, def:dc-ch6-2-4-principal, lem:dc-ch6-2-8-euclidean-opens, ex:dc-ch6-2-8}
  \leanok
  \dcref{ch6:2.4}
  An immersion $f:M^n\to\overline{M}^{n+1}$ has image a \emph{hypersurface}.
  For $p\in M$ and a unit normal $\eta\in(T_pM)^\perp$, the symmetric operator
  $S_\eta:T_pM\to T_pM$ has an orthonormal basis of
  eigenvectors $\{e_1,\dots,e_n\}$ of $T_pM$ with real eigenvalues
  $\lambda_1,\dots,\lambda_n$, $S_\eta(e_i)=\lambda_i e_i$. If $M$ and $\overline{M}$
  are both orientable and oriented, then $\eta$ is uniquely determined by requiring
  $\{e_1,\dots,e_n\}$ to be in the orientation of $M$ and $\{e_1,\dots,e_n,\eta\}$ in
  the orientation of $\overline{M}$. The $e_i$ are \emph{principal directions} and the
  $\lambda_i=k_i$ are \emph{principal curvatures} of $f$. The symmetric functions of the
  $\lambda_i$ are invariants of the immersion: $\det(S_\eta)=\lambda_1\cdots\lambda_n$
  is the \emph{Gauss-Kronecker curvature} and
  $\frac{1}{n}(\lambda_1+\cdots+\lambda_n)$ the \emph{mean curvature} of $f$.
  When $\overline{M}=\mathbb{R}^{n+1}$, a local unit normal $N$ defines the
  \emph{Gauss spherical map} $g:M^n\to S_1^n$ by $g(q)=N(q)$. Under the natural
  identification $T_qM=T_{g(q)}S_1^n$, its differential is
  $$
  dg_q(x)=\overline{\nabla}_xN=-S_{N(q)}(x).
  $$
\end{example}

\begin{remark}[compact hypersurfaces with non-vanishing Gauss-Kronecker curvature]
  \label{rem:dc-ch6-2-4-covering}
  \uses{ex:dc-ch6-2-4}
  \notready
  \notready
  If $M^n$, $n\geq2$, is connected, compact, and orientable and admits an
  immersion into $\mathbb R^{n+1}$ with nonzero Gauss--Kronecker curvature,
  then its Gauss map is a local diffeomorphism. Compactness makes it a covering
  of the simply connected sphere, hence $M\cong S^n$.
\end{remark}

\begin{definition}[principal curvatures and directions; Gauss-Kronecker and mean curvature]
  \label{def:dc-ch6-2-4-principal}
  \lean{Riemannian.DCImmersedPatch.shapeOperatorHom, Riemannian.DCImmersedPatch.shapeOperatorHom_isSymmetric, Riemannian.DCImmersedPatch.principalCurvatures, Riemannian.DCImmersedPatch.principalDirections, Riemannian.DCImmersedPatch.principalDirection, Riemannian.DCImmersedPatch.metricInner_principalDirection, Riemannian.DCImmersedPatch.shapeOperatorAt_principalDirection, Riemannian.DCImmersedPatch.secondFundScalarAt_principalDirection, Riemannian.DCImmersedPatch.gaussKroneckerCurvature, Riemannian.DCImmersedPatch.meanCurvature, Riemannian.DCImmersedPatch.trace_shapeOperatorHom, Riemannian.DCImmersedPatch.meanCurvature_eq_sum, Riemannian.DCImmersedPatch.gaussKroneckerCurvature_eq_prod}
  \uses{def:dc-ch6-2-2}
  \leanok
  Let $p\in M$ and $\eta\in(T_pM)^\perp$. Since $S_\eta:T_pM\to T_pM$ is linear
  and self-adjoint with respect to the metric, the spectral theorem yields an
  orthonormal basis of eigenvectors $\{e_1,\dots,e_n\}$ of $T_pM$ with real
  eigenvalues $\lambda_1,\dots,\lambda_n$, $S_\eta(e_i)=\lambda_ie_i$, and
  $H_\eta(e_i,e_j)=\lambda_i\delta_{ij}$. The $e_i$ are the \emph{principal
  directions} and the $\lambda_i=k_i$ the \emph{principal curvatures} of $f$ at
  $p$ along $\eta$. Their symmetric functions are invariants of the immersion:
  $\det(S_\eta)=\lambda_1\cdots\lambda_n$ is the \emph{Gauss-Kronecker
  curvature} and $\frac1n\operatorname{trace}(S_\eta)
  =\frac1n(\lambda_1+\cdots+\lambda_n)$ the \emph{mean curvature} of $f$ at $p$
  along $\eta$.
\end{definition}

\begin{theorem}[Gauss]
  \label{thm:dc-ch6-2-5}
  \lean{Riemannian.DCImmersedPatch.gauss_theorem, Riemannian.DCImmersedPatch.inducedCurvature_inner_sub_curvature_inner}
  \uses{prop:dc-ch6-2-1, def:dc-ch6-3-curvatures, def:dc-ch4-3-2, prop:dc-ch6-3-1}
  \leanok
  \dcref{ch6:2.5}
  Let $p\in M$ and let $x,y$ be orthonormal vectors in $T_pM$. Then

  $$
  K(x,y)-\overline{K}(x,y)=\langle B(x,x),B(y,y)\rangle-|B(x,y)|^2,\qquad(1)
  $$

  where $K(x,y)$ and $\overline{K}(x,y)$ denote the sectional curvatures of $M$ and
  $\overline{M}$, respectively, in the plane generated by $x$ and $y$.
\end{theorem}
\begin{proof}
  \leanok
  \sketch
  Extend $x,y$ to tangent fields $X,Y$ and choose an orthonormal normal frame
  $E_i$. Insert the Gauss decomposition in the curvature definition. The normal
  component of the bracket term pairs trivially with $Y$; pairing the two
  second-covariant-derivative terms with $Y$ gives respectively
  $-\sum_iH_i(X,X)H_i(Y,Y)$ and $-\sum_iH_i(X,Y)^2$, while the tangential
  terms reproduce the curvature of $M$. Their difference is (1). Equivalently,
  (1) is the $Z=X,T=Y$ specialization of the Gauss equation
  (\cref{prop:dc-ch6-3-1}) for an orthonormal pair.
\end{proof}

\begin{remark}[Gauss formula for a hypersurface]
  \label{rem:dc-ch6-2-6}
  \lean{Riemannian.DCImmersedPatch.secondFundFormAt_eq_secondFundScalarAt_smul, Riemannian.DCImmersedPatch.hypersurface_gauss}
  \uses{thm:dc-ch6-2-5, def:dc-ch6-2-4-principal}
  \leanok
  \dcref{ch6:2.6}
  In the case of a hypersurface $f:M^n\to\overline{M}^{n+1}$, let $p\in M$,
  $\eta\in(T_pM)^\perp$, $|\eta|=1$, and $\{e_1,\dots,e_n\}$ an orthonormal basis of $T_pM$ in
  which $S_\eta=S$ is diagonal, $S(e_i)=\lambda_ie_i$. Then $H(e_i,e_i)=\lambda_i$
  and $H(e_i,e_j)=0$ if $i\neq j$, so (1) becomes
  $$
  K(e_i,e_j)-\overline{K}(e_i,e_j)=\lambda_i\lambda_j.\qquad(4)
  $$
\end{remark}

\begin{lemma}[Koszul formula for the induced connection]
  \label{lem:dc-ch6-2-7-koszul}
  \lean{Riemannian.DCImmersedPatch.inducedCov_koszul}
  \uses{prop:dc-ch6-2-induced-connection}
  \leanok
  For fields $X,Y,Z$ tangent to the patch,
  $$
  2\langle Z,\nabla_YX\rangle
    = X\langle Y,Z\rangle + Y\langle Z,X\rangle - Z\langle X,Y\rangle
      - \langle[X,Z],Y\rangle - \langle[Y,Z],X\rangle - \langle[X,Y],Z\rangle.
  $$
  Every quantity on the right-hand side — inner products of tangent fields,
  their derivatives along tangent fields, and brackets of tangent fields — is
  intrinsic to the immersed manifold: the induced connection is determined by
  the induced metric.
\end{lemma}
\begin{proof}
  \leanok
  \sketch
  Write the compatibility identity of \cref{prop:dc-ch6-2-induced-connection}
  for the three cyclic orderings,
  $X\langle Y,Z\rangle=\langle\nabla_XY,Z\rangle+\langle Y,\nabla_XZ\rangle$,
  $Y\langle Z,X\rangle=\langle\nabla_YZ,X\rangle+\langle Z,\nabla_YX\rangle$,
  $Z\langle X,Y\rangle=\langle\nabla_ZX,Y\rangle+\langle X,\nabla_ZY\rangle$,
  and substitute the symmetry identities
  $\nabla_XY-\nabla_YX=[X,Y]$, $\nabla_XZ-\nabla_ZX=[X,Z]$,
  $\nabla_YZ-\nabla_ZY=[Y,Z]$ of the same proposition. Adding the first two
  identities and subtracting the third, all covariant-derivative terms cancel
  against bracket terms except $2\langle Z,\nabla_YX\rangle$, which gives the
  formula.
\end{proof}

\begin{lemma}[orthonormal pairs spanning a plane give equal curvature values]
  \label{lem:dc-ch6-2-7-plane}
  \lean{Riemannian.linearIndependent_pair_of_orthonormal, Riemannian.IsAlgCurvatureForm.apply_eq_of_orthonormal_of_mem_span, Riemannian.IsAlgCurvatureForm.apply_eq_of_orthonormal_of_finrank_eq_two}
  \uses{prop:dc-ch4-2-5, prop:dc-ch4-3-1}
  \leanok
  Let $B$ be an algebraic curvature form (\cref{prop:dc-ch4-2-5}) on an inner
  product space $V$, and let $\{u_1,u_2\}$ and $\{v_1,v_2\}$ be orthonormal
  pairs with $v_1,v_2\in\operatorname{span}\{u_1,u_2\}$ — in particular any two
  orthonormal pairs when $\dim V=2$. Then
  $$
  B(v_1,v_2,v_1,v_2)=B(u_1,u_2,u_1,u_2).
  $$
\end{lemma}
\begin{proof}
  \leanok
  \sketch
  Write $v_1=au_1+bu_2$ and $v_2=cu_1+du_2$. Orthonormality of the two pairs
  gives $a^2+b^2=1$, $c^2+d^2=1$ and $ac+bd=0$, so by the Lagrange identity
  $$
  (ad-bc)^2=(a^2+b^2)(c^2+d^2)-(ac+bd)^2=1,
  $$
  and the change of basis is invertible. An orthonormal pair is linearly
  independent (pair a vanishing combination against $u_1$ and $u_2$), so by
  the basis independence of the sectional curvature (\cref{prop:dc-ch4-3-1})
  $$
  \frac{B(v_1,v_2,v_1,v_2)}{|v_1\wedge v_2|^2}
    =\frac{B(u_1,u_2,u_1,u_2)}{|u_1\wedge u_2|^2},
  $$
  and both denominators equal $1$ by orthonormality. When $\dim V=2$ the pair
  $\{u_1,u_2\}$ is a basis, so the span hypothesis is automatic.
\end{proof}

\begin{definition}[pushforward of fields along a diffeomorphism of Euclidean opens]
  \label{def:dc-ch6-2-7-pushforward}
  \lean{Riemannian.opensMapExt, Riemannian.opensFDeriv, Riemannian.opensPushforward, Riemannian.fderiv_extendZero_comp, Riemannian.opensFDeriv_symm_comp}
  \uses{lem:dc-ch6-2-8-euclidean-opens}
  \leanok
  Let $S\subseteq F$ and $T\subseteq G$ be open subsets of Euclidean spaces and
  $\varphi:S\to T$ a smooth map. Its \emph{derivative} $d\varphi_p:F\to G$ at
  $p\in S$ is the derivative of any local extension of $\varphi$ to the
  ambient space; it obeys the chain rule
  $d(g\circ\varphi)_p=dg_{\varphi(p)}\circ d\varphi_p$, and if $\varphi$ is a
  diffeomorphism with inverse $\psi$ then
  $d\psi_{\varphi(p)}\circ d\varphi_p=\mathrm{id}$. For a diffeomorphism
  $\varphi$ and a smooth vector field $X$ on $S$, the \emph{pushforward}
  $$
  (\varphi_*X)(q)=d\varphi_{\varphi^{-1}(q)}\bigl(X(\varphi^{-1}(q))\bigr)
  $$
  is a smooth vector field on $T$ (the derivative of a smooth map depends
  smoothly on the base point).
\end{definition}

\begin{lemma}[naturality of derivations and brackets under pushforward]
  \label{lem:dc-ch6-2-7-pushforward-naturality}
  \lean{Riemannian.dir_opensPushforward, Riemannian.DCLieBracket_opensPushforward, Riemannian.bracketField_opensPushforward, Riemannian.metricInner_field_contMDiff_opens}
  \uses{def:dc-ch6-2-7-pushforward, lem:dc-ch0-5-2}
  \leanok
  Let $\varphi$ be a diffeomorphism of Euclidean opens. For smooth vector
  fields $X,Y$ and a smooth scalar $h$,
  $$
  (\varphi_*X)(h\circ\varphi^{-1})=(Xh)\circ\varphi^{-1}
  \qquad\text{and}\qquad
  [\varphi_*X,\varphi_*Y]=\varphi_*[X,Y].
  $$
\end{lemma}
\begin{proof}
  \leanok
  \sketch
  The first identity is the chain rule of \cref{def:dc-ch6-2-7-pushforward}
  together with $d\psi_{\varphi(p)}\circ d\varphi_p=\mathrm{id}$, where
  $\psi=\varphi^{-1}$. For the bracket, compute the derivative of the
  pushforward field $\varphi_*Y=(d\varphi\circ Y)\circ\psi$ at
  $q=\varphi(p)$ by the chain and product rules:
  $$
  d(\varphi_*Y)_q\bigl(d\varphi_p(X_p)\bigr)
    = d^2\varphi_p(X_p,Y_p)+d\varphi_p\bigl(dY_p(X_p)\bigr).
  $$
  Subtracting the same expression with $X$ and $Y$ interchanged, the
  second-derivative terms cancel by symmetry of $d^2\varphi_p$,
  leaving
  $[\varphi_*X,\varphi_*Y](q)
    =d\varphi_p\bigl(dY_p(X_p)-dX_p(Y_p)\bigr)=d\varphi_p([X,Y](p))$,
  the bracket of fields on a Euclidean open being the commutator of directional
  derivatives.
\end{proof}

\begin{definition}[isometry of immersed patches]
  \label{def:dc-ch6-2-7-isometry}
  \lean{Riemannian.DCPatchIsometry.isTangentField_opensPushforward, Riemannian.DCPatchIsometry.metricInner_opensPushforward}
  \uses{def:dc-ch6-2-patch, def:dc-ch1-2-2, def:dc-ch6-2-7-pushforward}
  \leanok
  Let $D$ and $D'$ be immersed patches in open subsets of Euclidean spaces. An
  \emph{isometry} $\varphi:D\to D'$ is a diffeomorphism of the ambient opens
  whose derivative carries each tangent space $T_pM$ into
  $T_{\varphi(p)}M'$ (and, for the inverse, $T_qM'$ into
  $T_{\varphi^{-1}(q)}M$) and preserves the inner product of tangent vectors:
  $$
  \langle d\varphi_p(v),d\varphi_p(w)\rangle=\langle v,w\rangle
  \qquad\text{for } v,w\in T_pM.
  $$
  No condition is imposed on $d\varphi$ transverse to the tangent distributions;
  tangent fields push forward to tangent fields and tangent inner products are
  preserved pointwise.
\end{definition}

\begin{lemma}[isometries preserve the induced connection]
  \label{lem:dc-ch6-2-7-connection}
  \lean{Riemannian.DCPatchIsometry.inducedCov_opensPushforward}
  \uses{def:dc-ch6-2-7-isometry, lem:dc-ch6-2-7-koszul, lem:dc-ch6-2-7-pushforward-naturality}
  \leanok
  Let $\varphi:D\to D'$ be an isometry of immersed patches in flat Euclidean
  ambients, and let $\nabla,\nabla'$ be the induced connections. Then for
  tangent fields $X,Y$,
  $$
  \nabla'_{\varphi_*X}(\varphi_*Y)=\varphi_*(\nabla_XY).
  $$
\end{lemma}
\begin{proof}
  \leanok
  \sketch
  Both sides are tangent fields, so by positive-definiteness it suffices that
  their inner products against every tangent vector at every $q=\varphi(p)$
  agree. A tangent vector $w\in T_qM'$ is $d\varphi_p(z)$ for the tangent
  vector $z=d\varphi^{-1}_q(w)\in T_pM$; extend $z$ to a tangent field $Z$.
  By the Koszul formula (\cref{lem:dc-ch6-2-7-koszul}) applied upstairs to
  the tangent fields $\varphi_*X,\varphi_*Y,\varphi_*Z$, each term of
  $2\langle\varphi_*Z,\nabla'_{\varphi_*X}\varphi_*Y\rangle$ is a
  derivative of an inner product of tangent fields along a tangent field, or
  an inner product with a bracket of tangent fields; by
  \cref{lem:dc-ch6-2-7-pushforward-naturality} and the metric preservation of
  \cref{def:dc-ch6-2-7-isometry}, each equals the corresponding term of the
  Koszul formula downstairs, so
  $\langle\varphi_*Z,\nabla'_{\varphi_*X}\varphi_*Y\rangle(q)
    =\langle Z,\nabla_XY\rangle(p)
    =\langle\varphi_*Z,\varphi_*(\nabla_XY)\rangle(q)$,
  the last step again by metric preservation, since $\nabla_XY$ is tangent.
\end{proof}

\begin{lemma}[isometries preserve the induced curvature]
  \label{lem:dc-ch6-2-7-curvature}
  \lean{Riemannian.DCPatchIsometry.inducedCurvature_opensPushforward, Riemannian.DCPatchIsometry.inducedCurvature_inner_opensPushforward}
  \uses{lem:dc-ch6-2-7-connection, def:dc-ch6-3-curvatures}
  \leanok
  Under the hypotheses of \cref{lem:dc-ch6-2-7-connection}, for tangent fields
  $X,Y,Z$,
  $$
  R'(\varphi_*X,\varphi_*Y)(\varphi_*Z)=\varphi_*\bigl(R(X,Y)Z\bigr),
  $$
  and consequently the sectional numerator is invariant:
  $\langle R'(\varphi_*X,\varphi_*Y)\varphi_*X,\varphi_*Y\rangle(\varphi(p))
    =\langle R(X,Y)X,Y\rangle(p)$.
\end{lemma}
\begin{proof}
  \leanok
  \sketch
  Expand $R(X,Y)Z=\nabla_Y\nabla_XZ-\nabla_X\nabla_YZ+\nabla_{[X,Y]}Z$
  (\cref{def:dc-ch6-3-curvatures}) and apply
  \cref{lem:dc-ch6-2-7-connection} to each covariant derivative — the
  intermediate fields $\nabla_XZ$, $\nabla_YZ$ and the bracket $[X,Y]$ are
  again tangent, the bracket pushes forward by
  \cref{lem:dc-ch6-2-7-pushforward-naturality}, and the pushforward is
  pointwise linear. The numerator identity follows by pairing with
  $\varphi_*Y$ and metric preservation, $R(X,Y)X$ being tangent.
\end{proof}

\begin{lemma}[the pointwise Gauss form of a patch in a flat ambient]
  \label{lem:dc-ch6-2-7-gauss-form}
  \lean{Riemannian.DCImmersedPatch.gaussFormAt, Riemannian.DCImmersedPatch.isAlgCurvatureForm_gaussFormAt, Riemannian.inducedCurvature_inner_eq_gaussFormAt}
  \uses{def:dc-ch6-2-2, prop:dc-ch6-3-1, prop:dc-ch4-2-5, lem:dc-ch6-2-8-euclidean-opens}
  \leanok
  For an immersed patch and $p\in M$, the \emph{Gauss form}
  $$
  A(x,y,z,t)=\langle B(y,t),B(x,z)\rangle-\langle B(x,t),B(y,z)\rangle,
  \qquad x,y,z,t\in T_pM,
  $$
  built from the pointwise second fundamental form of \cref{def:dc-ch6-2-2},
  is an algebraic curvature form on $T_pM$ (\cref{prop:dc-ch4-2-5}). If the
  ambient space is a flat Euclidean open, then for tangent fields
  $$
  \langle R(X,Y)Z,T\rangle(p)=A\bigl(X_p,Y_p,Z_p,T_p\bigr):
  $$
  the curvature pairing depends only on the values of the fields at $p$.
\end{lemma}
\begin{proof}
  \uses{prop:dc-ch6-2-1}
  \leanok
  \sketch
  Multilinearity of $A$ follows from bilinearity of the pointwise $B$ and of
  the inner product. Antisymmetry in $(x,y)$ is immediate from the shape of
  $A$; antisymmetry in $(z,t)$ follows using the symmetry of the inner
  product; the first Bianchi identity follows from the symmetry
  $B(x,y)=B(y,x)$ on tangent vectors (\cref{prop:dc-ch6-2-1} read at a
  point): the six terms of the cyclic sum cancel in pairs. The identity with
  the curvature pairing is the Gauss equation (\cref{prop:dc-ch6-3-1}) with
  flat ambient curvature $\overline{R}\equiv 0$
  (\cref{lem:dc-ch6-2-8-euclidean-opens}), with $B$ evaluated pointwise as in
  \cref{def:dc-ch6-2-2}.
\end{proof}

\begin{remark}[isometry invariance of Gaussian curvature]
  \label{rem:dc-ch6-2-7}
  \lean{Riemannian.DCPatchIsometry.theorema_egregium}
  \uses{rem:dc-ch6-2-6, def:dc-ch6-2-7-isometry, def:dc-ch6-2-4-principal}
  \leanok
  \dcref{ch6:2.7}
  In the case $M=M^2\subset\overline{M}=\mathbb{R}^3$, the product $\lambda_1\lambda_2$
  of the principal curvatures coincides with the Gaussian curvature of the surface.
  \cref{rem:dc-ch6-2-6} then shows that the Gaussian curvature coincides with the sectional
  curvature of the surface, and gives the \emph{Theorema Egregium}:
  the Gaussian curvature of $M^2\subset\mathbb{R}^3$ is an invariant under
  isometries. Precisely: let $D,D'$ be two-dimensional hypersurface patches in
  flat Euclidean ambients, with unit normals $\eta$ at $p$ and $\eta'$ at
  $\varphi(p)$ spanning the normal lines, and let $\varphi:D\to D'$ be an
  isometry (\cref{def:dc-ch6-2-7-isometry}). Then
  $$
  \lambda_1(p)\,\lambda_2(p)=\lambda'_1(\varphi(p))\,\lambda'_2(\varphi(p)).
  $$
\end{remark}
\begin{proof}
  \uses{lem:dc-ch6-2-7-curvature, lem:dc-ch6-2-7-gauss-form,
    lem:dc-ch6-2-7-plane}
  \leanok
  \sketch
  Let $e_1,e_2$ be principal directions at $p$, extended to tangent fields
  $E_1,E_2$, and $e'_1,e'_2$ principal directions at $\varphi(p)$, extended to
  $E'_1,E'_2$. Since the ambients are flat, the Gauss formula for a
  hypersurface (\cref{rem:dc-ch6-2-6}) reads
  $$
  \lambda_1\lambda_2=\langle R(E_1,E_2)E_1,E_2\rangle(p),
  \qquad
  \lambda'_1\lambda'_2=\langle R'(E'_1,E'_2)E'_1,E'_2\rangle(\varphi(p)).
  $$
  By the curvature invariance (\cref{lem:dc-ch6-2-7-curvature}),
  $\langle R(E_1,E_2)E_1,E_2\rangle(p)
    =\langle R'(\varphi_*E_1,\varphi_*E_2)\varphi_*E_1,\varphi_*E_2\rangle(\varphi(p))$.
  By \cref{lem:dc-ch6-2-7-gauss-form} both curvature pairings at
  $\varphi(p)$ are values of the Gauss form $A$ of $D'$: on the pair
  $u_i=d\varphi_p(e_i)$, and on the pair $e'_i$, respectively. Both pairs are
  orthonormal in the two-dimensional plane $T_{\varphi(p)}M'$ — the first
  because $\varphi$ is an isometry — so the two values of $A$ agree by
  \cref{lem:dc-ch6-2-7-plane}, and the chain
  $\lambda_1\lambda_2
    =A(u_1,u_2,u_1,u_2)=A(e'_1,e'_2,e'_1,e'_2)=\lambda'_1\lambda'_2$
  closes.
\end{proof}

\begin{lemma}[the Euclidean ambient: flat connection, vanishing curvature]
  \label{lem:dc-ch6-2-8-euclidean}
  \lean{Riemannian.SmoothVectorField.contDiff, Riemannian.DCLieBracket_eq_fderiv, Riemannian.euclideanCov, Riemannian.euclideanConnection, Riemannian.euclideanConnection_isLeviCivita, Riemannian.euclideanConnection_curvature, Riemannian.euclideanConnection_isConstantCurvature_zero}
  \uses{def:dc-ch1-2-1, thm:dc-ch2-3-6, def:dc-ch4-2-1, def:dc-ch6-3-constant-curvature}
  \leanok
  On Euclidean space — a real inner-product space $F$ with the metric of
  \cref{def:dc-ch1-2-1} given by the ambient inner product — the ordinary
  directional derivative $\nabla_XY=dY(X)$ is an affine connection, it is the
  Levi-Civita connection of the Euclidean metric, and its curvature tensor
  vanishes identically; in particular Euclidean space has constant sectional
  curvature $0$ in the field form of \cref{def:dc-ch6-3-constant-curvature}.
\end{lemma}
\begin{proof}
  \leanok
  \sketch
  On the vector-space model the tangent bundle is canonically trivial, so a
  smooth vector field \emph{is} a smooth map $F\to F$, and the manifold Lie
  bracket is the classical commutator $[X,Y]=dY(X)-dX(Y)$. Additivity and
  $\mathcal{D}$-homogeneity of $\nabla_XY=dY(X)$ in $X$ are linearity of
  $dY_p$; additivity and the Leibniz rule in $Y$ are the sum and product rules
  for the derivative. Symmetry is $dY(X)-dX(Y)=[X,Y]$; metric compatibility is
  the product rule for the inner product. For the curvature,
  $$
  R(X,Y)Z=\nabla_Y\nabla_XZ-\nabla_X\nabla_YZ+\nabla_{[X,Y]}Z ,
  $$
  the second-order terms cancel by equality of mixed derivatives, and the remaining first-order terms cancel exactly against
  $\nabla_{[X,Y]}Z=dZ([X,Y])$. Constant curvature $0$ is then immediate.
\end{proof}

\begin{lemma}[open subsets of Euclidean space as ambient manifolds]
  \label{lem:dc-ch6-2-8-euclidean-opens}
  \lean{Riemannian.opensEuclideanMetric, Riemannian.opensEuclideanConnection, Riemannian.opensEuclideanConnection_isLeviCivita, Riemannian.opensEuclideanConnection_curvature, Riemannian.opensEuclideanConnection_isConstantCurvature_zero, Riemannian.DCLieBracket_opens_eq_fderiv}
  \uses{def:dc-ch1-2-1, thm:dc-ch2-3-6, def:dc-ch4-2-1, def:dc-ch6-3-constant-curvature, lem:dc-ch6-2-8-euclidean}
  \leanok
  Let $A\subseteq F$ be an open subset of a Euclidean space, regarded as an
  open submanifold. A vector field on $A$ is canonically a smooth map
  $A\to F$, and under this identification $A$ carries the restricted Euclidean
  metric $\langle v,w\rangle_q=\langle v,w\rangle$, the Lie bracket is the
  commutator $[X,Y]=dY(X)-dX(Y)$, and the ordinary directional derivative
  $\nabla_XY=dY(X)$ is the Levi-Civita connection of this metric, with
  identically vanishing curvature. Thus $A$ is a flat ambient manifold of
  constant sectional curvature $0$ in the field form of
  \cref{def:dc-ch6-3-constant-curvature}.
\end{lemma}
\begin{proof}
  \leanok
  \sketch
  Smoothness of a field on $A$ and its derivative at $q\in A$ depend only on
  the germ at $q$, and $A$ is open in $F$, so both may be computed after
  extending the field to $F$. The pointwise arguments of
  \cref{lem:dc-ch6-2-8-euclidean} then apply verbatim at each $q\in A$:
  linearity and the Leibniz rule are the sum and product rules of the
  derivative, symmetry of the connection is the commutator identity, metric
  compatibility is the product rule for the inner product, and the curvature
  vanishes by equality of mixed derivatives.
\end{proof}

\begin{example}[curvature of $S^n$]
  \label{ex:dc-ch6-2-8}
  \lean{Riemannian.punctured, Riemannian.sphereTang, Riemannian.spherePatch, Riemannian.spherePatch_secondFundForm_apply, Riemannian.spherePatch_secondFundScalarAt, Riemannian.spherePatch_shapeOperatorAt, Riemannian.spherePatch_inducedCurvature_inner, Riemannian.sphere_sectionalCurvature_one}
  \uses{def:dc-ch6-2-patch, def:dc-ch6-2-B, def:dc-ch6-2-2, thm:dc-ch6-2-5, rem:dc-ch6-2-6, lem:dc-ch6-2-8-euclidean-opens}
  \leanok
  \dcref{ch6:2.8}
  The sphere of radius $r$ in Euclidean space has constant sectional curvature
  $1/r^2$; in particular, the unit sphere has constant sectional curvature $1$.
\end{example}
\begin{proof}
  \sketch
  With the inward unit normal $\eta=-p/r$, the second fundamental form is
  $H_\eta(x,y)=\langle x,y\rangle/r$, so $S_\eta=r^{-1}\mathrm{id}$. Gauss'
  equation against the flat ambient gives sectional curvature $1/r^2$.
\end{proof}

\begin{definition}[geodesic at a point; totally geodesic immersion]
  \label{def:dc-ch6-2-8-geodesic}
  \lean{Riemannian.DCImmersedPatch.IsGeodesicAt, Riemannian.DCImmersedPatch.IsTotallyGeodesic, Riemannian.DCImmersedPatch.isGeodesicAt_iff_secondFundFormAt_eq_zero}
  \source{docarmo:page-0147}
  \uses{def:dc-ch6-2-2}
  \leanok
  An immersion $f:M\to\overline{M}$ is \emph{geodesic at} $p\in M$ if for every
  $\eta\in(T_pM)^\perp$ the second fundamental form $\mathit{II}_\eta$ is
  identically zero at $p$ — equivalently, by polarization, $B$ vanishes at $p$.
  An immersion is \emph{totally geodesic} if it is geodesic at every
  $p\in M$.
\end{definition}

\begin{definition}[geodesic field at a point]
  \label{def:dc-ch6-2-9-geodesic-field}
  \lean{Riemannian.DCImmersedPatch.IsGeodesicFieldAt}
  \uses{def:dc-ch6-2-induced-connection, def:dc-ch6-2-8-geodesic}
  \leanok
  A tangent field $X\in\mathcal{X}(U)$ is a \emph{geodesic field at} $p$ when
  $(\nabla_X X)(p)=0$. The velocity field of a geodesic satisfies this condition
  along the geodesic.
\end{definition}

\begin{lemma}[existence of geodesic fields; first-order scalar data]
  \label{lem:dc-ch6-2-9-existence}
  \lean{Riemannian.exists_contMDiff_dirTangent_eq_one, Riemannian.AffineConnection.cov_apply_eq_zero_of_apply_eq_zero_left, Riemannian.DCImmersedPatch.exists_isGeodesicFieldAt}
  \uses{def:dc-ch6-2-9-geodesic-field, def:dc-ch6-2-induced-connection, lem:dc-ch6-2-proj}
  \leanok
  (a) For every $p\in\overline{M}$ and $v\in T_p\overline{M}$, $v\neq0$, there is a
  globally smooth $f:\overline{M}\to\mathbb{R}$ with $f(p)=0$ and $df_p(v)=1$.
  (b) Through every $x\in T_pM$ there passes a geodesic field at $p$: a tangent
  field $X$ with $X(p)=x$ and $(\nabla_X X)(p)=0$.
\end{lemma}
\begin{proof}
  \leanok
  \sketch
  (a) Choose by Hahn--Banach a continuous linear functional $\ell$ on the model
  space with $\ell(v)=1$, and let $\varphi$ be the coordinate reading of $\ell$
  in the chart at $p$, shifted so that $\varphi(p)=0$. Cutting $\varphi$ off by
  a smooth bump function equal to $1$ near $p$ and supported in the chart source
  yields a global smooth $f$ with $f(p)=0$; since the cut-off is constant near
  $p$ it does not affect the differential, and $df_p(v)=\ell(d(\mathrm{chart})_p
  v)=\ell(v)=1$ because the differential of the chart at its base point is the
  identity.

  (b) If $x=0$ the zero field works, by tensoriality of the covariant derivative
  in its direction slot. Otherwise let $X_0$ be a tangent extension of $x$
  (\cref{lem:dc-ch6-2-proj}), set $a:=(\nabla_{X_0}X_0)(p)\in T_pM$, let $W$ be a
  tangent extension of $-a$, and take $f$ as in (a) with $df_p(x)=1$, $f(p)=0$.
  Put $X:=X_0+fW$; then $X$ is tangent and $X(p)=x$. Expanding
  $\nabla_{X}X$ by additivity in both slots, $\mathcal{D}$-homogeneity in the
  direction slot and the Leibniz rule in the derived slot,
  $$
  \nabla_X X=\nabla_{X_0}X_0+f\,\nabla_{X_0}W+X_0(f)\,W+f\,\nabla_{W}X ,
  $$
  and at $p$ all terms carrying the factor $f(p)=0$ vanish, leaving
  $(\nabla_XX)(p)=a+X_0(f)(p)\,W(p)=a+1\cdot(-a)=0$.
\end{proof}

\begin{proposition}[totally geodesic criterion]
  \label{prop:dc-ch6-2-9}
  \source{docarmo:page-0147}
  \uses{def:dc-ch6-2-8-geodesic, def:dc-ch6-2-9-geodesic-field, lem:dc-ch6-2-9-existence, def:dc-ch6-2-B, def:dc-ch6-2-2, lem:dc-ch6-2-locality, prop:dc-ch6-2-1, lem:dc-ch6-2-9-field}
  \notready
  \dcref{ch6:2.9}
  An immersion $f:M\to\overline{M}$ is geodesic at $p\in M$ if and only if every
  geodesic of $M$ at $p$ is a geodesic of $\overline{M}$ at $p$; precisely: if
  and only if every geodesic field $X$ at $p$
  (\cref{def:dc-ch6-2-9-geodesic-field}) satisfies
  $(\overline{\nabla}_X X)(p)=0$. The two readings agree through the displayed
  identity of the proof: for any tangent field $X$ with $X(p)=x$ and any
  $\eta\in(T_pM)^\perp$,
  $$
  H_\eta(x,x)=\langle\overline{\nabla}_X X,\eta\rangle(p).
  $$
\end{proposition}
\begin{proof}
  \sketch
  For any tangent field $X$ with $X(p)=x$ and $\eta\in(T_pM)^\perp$, the Gauss
  decomposition $\overline{\nabla}_X X=\nabla_X X+B(X,X)$ paired against $\eta$
  gives
  $$
  H_\eta(x,x)=\langle B(x,x),\eta\rangle=\langle\overline{\nabla}_X X,\eta\rangle(p),
  $$
  since the tangential part $\nabla_X X$ pairs to zero against $\eta$.

  ($\Rightarrow$) If $f$ is geodesic at $p$ then $B$ vanishes at $p$, so for a
  geodesic field $X$ at $p$ the ambient acceleration
  $(\overline{\nabla}_X X)(p)=(\nabla_X X)(p)+B(x,x)=0+0=0$.

  ($\Leftarrow$) Let $\eta\in(T_pM)^\perp$ and $x\in T_pM$. By
  \cref{lem:dc-ch6-2-9-existence} there is a geodesic field $X$ at $p$ with
  $X(p)=x$; by hypothesis $(\overline{\nabla}_X X)(p)=0$, hence
  $H_\eta(x,x)=\langle\overline{\nabla}_X X,\eta\rangle(p)=0$. Thus every
  $\mathit{II}_\eta$ vanishes at $p$, i.e.\ $f$ is geodesic at $p$.
\end{proof}

\begin{remark}[geometric interpretation of sectional curvature]
  \label{rem:dc-ch6-2-9-interpretation}
  \lean{Riemannian.DCImmersedPatch.secondFundForm_apply_eq_zero_of_isGeodesicAt, Riemannian.DCImmersedPatch.inducedCurvature_inner_eq_curvature_inner_of_isGeodesicAt, Riemannian.DCImmersedPatch.sectionalCurvature_eq_of_isGeodesicAt, Riemannian.DCImmersedPatch.inducedCurvatureAt_inner_eq_of_isGeodesicAt}
  \source{docarmo:page-0147,docarmo:page-0148}
  \uses{prop:dc-ch6-2-9, thm:dc-ch6-2-5, def:dc-ch6-2-8-geodesic, def:dc-ch4-3-2}
  \leanok
  If $B\subset T_pM$ is a normal neighborhood and $\sigma\subset T_pM$ is
  two-dimensional, then $S=\exp_p(\sigma\cap B)$ is geodesic at $p$ and
  $$
  K_S(p)=K(p,\sigma).
  $$
  More generally, any submanifold geodesic at $p$ has the same sectional
  curvature as the ambient manifold on each tangent plane at $p$.
\end{remark}
\begin{proof}
  \sketch
  The exponential construction makes the second fundamental form of $S$ vanish
  at $p$. The equality then follows from Gauss' equation; the same argument
  applies to any submanifold geodesic at $p$.
\end{proof}

\begin{remark}[geodesic surfaces and constant sectional curvature]
  \label{rem:dc-ch6-2-9-exp-surface}
  \source{docarmo:page-0148}
  \uses{prop:dc-ch6-2-9}
  \notready
  \notready
  When $\exp_p$ is a smooth local diffeomorphism and
  $\exp_p|_{\sigma\cap B}$ is an immersion, the surface
  $S=\exp_p(\sigma\cap B)$ is geodesic at $p$. Linear subspaces and translates
  in Euclidean space, and intersections of linear subspaces with $S^n$, are
  totally geodesic. Conversely, if every two-plane in every tangent space is
  tangent to a totally geodesic submanifold, then sectional curvature is
  constant.
\end{remark}

\begin{definition}[minimal immersion; mean curvature vector]
  \label{def:dc-ch6-2-10}
  \lean{Riemannian.DCImmersedPatch.IsMinimal, Riemannian.DCImmersedPatch.meanCurvatureVector, Riemannian.DCImmersedPatch.meanCurvatureVector_mem, Riemannian.DCImmersedPatch.inner_meanCurvatureVector, Riemannian.DCImmersedPatch.meanCurvatureVector_unique, Riemannian.DCImmersedPatch.existsUnique_meanCurvatureVector, Riemannian.DCImmersedPatch.isMinimal_iff_meanCurvatureVector_eq_zero, Riemannian.DCImmersedPatch.shapeOperatorAt_add_normal, Riemannian.DCImmersedPatch.shapeOperatorAt_smul_normal}
  \source{docarmo:page-0148}
  \uses{def:dc-ch6-2-4-principal, def:dc-ch6-2-2}
  \leanok
  \dcref{ch6:2.10}
  An immersion $f:M\to\overline{M}$ is called \emph{minimal} if for every $p\in M$ and
  every $\eta\in(T_pM)^\perp$ the trace of $S_\eta$ is zero. Choosing an orthonormal frame
  $E_1,\dots,E_m$ of vectors in $\mathcal{X}(U)^\perp$, where $U$ is a neighborhood
  of $p$ in which $f$ is an embedding, we write at $p$
  $$
  B(x,y)=\sum_i H_i(x,y)E_i,\quad x,y\in T_pM,\ i=1,\dots,m,
  $$
  with $H_i=H_{E_i}$. The normal vector
  $$
  H=\frac{1}{n}\sum_i(\operatorname{trace}S_i)E_i,
  $$
  where $S_i=S_{E_i}$, does not depend on the chosen frame; it is the \emph{mean
  curvature vector} of $f$. Thus $f$ is minimal iff $H(p)=0$ for all $p\in M$.
\end{definition}

\section{The fundamental equations}

Use Latin letters $X,Y,Z$ for tangent fields and Greek letters $\xi,\eta,\zeta$
for normal fields. The tangential and normal components of
$\overline{\nabla}_X\eta$ are respectively $-S_\eta(X)$ and $\nabla_X^\perp\eta$:

$$
\nabla_X^\perp\eta=(\overline{\nabla}_X\eta)^N=\overline{\nabla}_X\eta-(\overline{\nabla}_X\eta)^T=\overline{\nabla}_X\eta+S_\eta(X).\qquad(5)
$$

\begin{definition}[shape operator and normal connection: the Weingarten decomposition]
  \label{def:dc-ch6-3-normal-connection}
  \lean{Riemannian.DCImmersedPatch.shapeOperator, Riemannian.DCImmersedPatch.normalCov, Riemannian.DCImmersedPatch.cov_eq_neg_shapeOperator_add_normalCov, Riemannian.DCImmersedPatch.shapeOperator_mem, Riemannian.DCImmersedPatch.normalCov_mem}
  \uses{def:dc-ch6-2-patch, def:dc-ch6-2-B}
  \leanok
  For fields $X$, $\eta$ on the patch set
  $S_\eta(X)=-(\overline{\nabla}_X\eta)^T$ (the \emph{shape operator}, tangent
  valued) and $\nabla^\perp_X\eta=(\overline{\nabla}_X\eta)^N$ (the \emph{normal
  connection}, normal valued), so that equation (5) reads

  $$
  \overline{\nabla}_X\eta=-S_\eta(X)+\nabla^\perp_X\eta.
  $$
\end{definition}

\begin{proposition}[connection laws of the normal connection]
  \label{prop:dc-ch6-3-normal-connection-laws}
  \lean{Riemannian.DCImmersedPatch.normalCov_add_left, Riemannian.DCImmersedPatch.normalCov_smul_left, Riemannian.DCImmersedPatch.normalCov_add_right, Riemannian.DCImmersedPatch.normalCov_smul_right}
  \uses{def:dc-ch6-3-normal-connection, lem:dc-ch6-2-proj}
  \leanok
  The normal connection has the usual properties of a connection: it is
  additive in both arguments, $\mathcal{D}$-homogeneous in the direction
  ($\nabla^\perp_{fX}\eta=f\nabla^\perp_X\eta$), and satisfies the Leibniz rule
  $\nabla^\perp_X(f\eta)=f\nabla^\perp_X\eta+X(f)\,\eta$ for $\eta$ normal.
\end{proposition}
\begin{proof}
  \leanok
  \sketch
  Immediate from the corresponding laws of $\overline{\nabla}$ and the
  additivity and $\mathcal{D}$-homogeneity of the normal projection
  (\cref{lem:dc-ch6-2-proj}); in the Leibniz rule the first-order term
  $X(f)\,\eta$ is already normal because $\eta$ is.
\end{proof}

\begin{lemma}[Weingarten pairing and self-adjointness]
  \label{lem:dc-ch6-3-weingarten}
  \lean{Riemannian.DCImmersedPatch.inner_shapeOperator_apply, Riemannian.DCImmersedPatch.inner_shapeOperator_symm}
  \uses{def:dc-ch6-3-normal-connection, prop:dc-ch6-2-1, def:dc-ch6-2-B}
  \leanok
  For $Y$ tangent and $\eta$ normal,

  $$
  \langle S_\eta(X),Y\rangle=\langle B(X,Y),\eta\rangle
  $$

  at every point. Consequently, for tangent $X,Y$,
  $\langle S_\eta(X),Y\rangle=\langle X,S_\eta(Y)\rangle$: the shape operator
  is self-adjoint with respect to the metric.
\end{lemma}
\begin{proof}
  \leanok
  \sketch
  Since $\langle\eta,Y\rangle=0$ identically, compatibility of
  $\overline{\nabla}$ with the metric gives
  $0=X\langle\eta,Y\rangle
  =\langle\overline{\nabla}_X\eta,Y\rangle+\langle\eta,\overline{\nabla}_XY\rangle$.
  In the first term only the tangential part of $\overline{\nabla}_X\eta$
  survives against the tangent field $Y$, giving $-\langle S_\eta(X),Y\rangle$;
  in the second only the normal part of $\overline{\nabla}_XY$ survives against
  the normal field $\eta$, giving $\langle B(X,Y),\eta\rangle$. The
  self-adjointness follows from the symmetry of $B$
  (\cref{prop:dc-ch6-2-1}):
  $\langle S_\eta(X),Y\rangle=\langle B(X,Y),\eta\rangle
  =\langle B(Y,X),\eta\rangle=\langle S_\eta(Y),X\rangle$.
\end{proof}

\begin{definition}[curvatures of the induced and normal connections]
  \label{def:dc-ch6-3-curvatures}
  \lean{Riemannian.DCImmersedPatch.inducedCurvature, Riemannian.DCImmersedPatch.inducedCurvature_mem, Riemannian.DCImmersedPatch.normalCurvature, Riemannian.DCImmersedPatch.normalCurvature_mem}
  \uses{def:dc-ch6-2-induced-connection, def:dc-ch6-3-normal-connection, def:dc-ch4-2-1}
  \leanok
  The \emph{curvature of the induced connection} is
  $R(X,Y)Z=\nabla_Y\nabla_XZ-\nabla_X\nabla_YZ+\nabla_{[X,Y]}Z$ (tangent
  valued), and the \emph{normal curvature} is
  $R^\perp(X,Y)\eta=\nabla^\perp_Y\nabla^\perp_X\eta
  -\nabla^\perp_X\nabla^\perp_Y\eta+\nabla^\perp_{[X,Y]}\eta$ (normal valued);
  on tangent fields $R$ is the curvature of the immersed manifold.
\end{definition}

\begin{lemma}[decomposition of the ambient curvature]
  \label{lem:dc-ch6-3-decomposition}
  \lean{Riemannian.DCImmersedPatch.curvature_decomposition, Riemannian.DCImmersedPatch.curvature_normal_decomposition}
  \uses{def:dc-ch6-3-curvatures, lem:dc-ch6-2-proj, def:dc-ch6-2-B, def:dc-ch6-3-normal-connection}
  \leanok
  For any fields $X,Y,Z,\eta$ on the patch,

  $$
  \overline{R}(X,Y)Z=R(X,Y)Z+B(Y,\nabla_X Z)-B(X,\nabla_Y Z)+B([X,Y],Z)-S_{B(X,Z)}Y+S_{B(Y,Z)}X+\nabla_Y^\perp B(X,Z)-\nabla_X^\perp B(Y,Z)\qquad(6)
  $$

  and

  $$
  \overline{R}(X,Y)\eta=R^\perp(X,Y)\eta-S_{\nabla_X^\perp\eta}Y+S_{\nabla_Y^\perp\eta}X-\nabla_Y(S_\eta X)+\nabla_X(S_\eta Y)-B(Y,S_\eta X)+B(X,S_\eta Y)-S_\eta([X,Y]).
  $$
\end{lemma}
\begin{proof}
  \leanok
  \sketch
  Substitute the Gauss decomposition
  $\overline{\nabla}_XW=\nabla_XW+B(X,W)$ (\cref{def:dc-ch6-2-B}) and the
  Weingarten decomposition
  $\overline{\nabla}_X\eta=-S_\eta(X)+\nabla^\perp_X\eta$
  (\cref{def:dc-ch6-3-normal-connection}) into
  $\overline{R}(X,Y)W=\overline{\nabla}_Y\overline{\nabla}_XW
  -\overline{\nabla}_X\overline{\nabla}_YW+\overline{\nabla}_{[X,Y]}W$ and
  collect terms: for the first identity split
  $\overline{\nabla}_XZ=\nabla_XZ+B(X,Z)$, then split
  $\overline{\nabla}_Y(\nabla_XZ)=\nabla_Y\nabla_XZ+B(Y,\nabla_XZ)$ and
  $\overline{\nabla}_Y(B(X,Z))=-S_{B(X,Z)}Y+\nabla^\perp_YB(X,Z)$, and likewise
  for the terms with $X$ and $Y$ interchanged and for
  $\overline{\nabla}_{[X,Y]}Z$; for the second identity split
  $\overline{\nabla}_X\eta=-S_\eta(X)+\nabla^\perp_X\eta$, then
  $\overline{\nabla}_Y(S_\eta X)=\nabla_Y(S_\eta X)+B(Y,S_\eta X)$ and
  $\overline{\nabla}_Y(\nabla^\perp_X\eta)
  =-S_{\nabla^\perp_X\eta}Y+\nabla^\perp_Y\nabla^\perp_X\eta$, and likewise for
  the remaining terms.
\end{proof}

\begin{proposition}[Gauss equation and Ricci equation]
  \label{prop:dc-ch6-3-1}
  \lean{Riemannian.DCImmersedPatch.gauss_equation, Riemannian.DCImmersedPatch.ricci_equation}
  \uses{def:dc-ch6-3-curvatures, lem:dc-ch6-3-decomposition, lem:dc-ch6-3-weingarten}
  \leanok
  \dcref{ch6:3.1}
  The equations are:
  (a) \emph{Gauss equation}
  $$
  \langle\overline{R}(X,Y)Z,T\rangle=\langle R(X,Y)Z,T\rangle-\langle B(Y,T),B(X,Z)\rangle+\langle B(X,T),B(Y,Z)\rangle.
  $$
  (b) \emph{Ricci equation}
  $$
  \langle\overline{R}(X,Y)\eta,\zeta\rangle-\langle R^\perp(X,Y)\eta,\zeta\rangle=\langle[S_\eta,S_\zeta]X,Y\rangle,
  $$
  where $[S_\eta,S_\zeta]$ denotes the operator $S_\eta\circ S_\zeta-S_\zeta\circ S_\eta$.
\end{proposition}
\begin{proof}
  \leanok
  \sketch
  Substitute the Gauss and Weingarten decompositions into the ambient curvature
  operator and collect tangent and normal parts, as in
  \cref{lem:dc-ch6-3-decomposition}. Pairing the first resulting identity with
  $T$ removes normal terms; the Weingarten pairing converts the remaining shape
  terms to $B$-pairings and yields Gauss' equation. Pairing the second identity
  with $\zeta$ removes tangent terms and gives
  $\langle R^\perp(X,Y)\eta,\zeta\rangle$ plus the commutator term
  $\langle[S_\eta,S_\zeta]X,Y\rangle$, which is Ricci's equation.
\end{proof}

\begin{remark}[Gauss equation as a sectional-curvature identity]
  \label{rem:dc-ch6-3-2}
  \lean{Riemannian.DCImmersedPatch.inducedCurvature_inner_sub_curvature_inner}
  \uses{thm:dc-ch6-2-5, prop:dc-ch6-3-1}
  \leanok
  \dcref{ch6:3.2}
  \cref{thm:dc-ch6-2-5} is a special case of Gauss' equation:
  taking $Z=X$ and $T=Y$ in the Gauss equation and using the symmetry of $B$
  yields $\langle R(X,Y)X,Y\rangle-\langle\overline{R}(X,Y)X,Y\rangle
  =\langle B(X,X),B(Y,Y)\rangle-|B(X,Y)|^2$.
\end{remark}

\begin{definition}[constant sectional curvature, field form]
  \label{def:dc-ch6-3-constant-curvature}
  \lean{Riemannian.AffineConnection.IsConstantCurvature}
  \uses{def:dc-ch4-2-1, lem:dc-ch4-3-4}
  \leanok
  A Riemannian manifold with Riemannian connection $\overline{\nabla}$ has
  \emph{constant sectional curvature} $K_0$ when

  $$
  \langle\overline{R}(X,Y)Z,W\rangle=K_0\big(\langle X,Z\rangle\langle Y,W\rangle-\langle Y,Z\rangle\langle X,W\rangle\big)
  $$

  for all fields $X,Y,Z,W$ — the characterization of constant sectional
  curvature of \cref{lem:dc-ch4-3-4}.
\end{definition}

\begin{remark}[flat normal bundle]
  \label{rem:dc-ch6-3-3}
  \lean{Riemannian.DCImmersedPatch.ricci_equation_of_isConstantCurvature, Riemannian.DCImmersedPatch.HasFlatNormalBundle, Riemannian.DCImmersedPatch.ShapeOperatorsCommuteAt, Riemannian.DCImmersedPatch.hasFlatNormalBundle_iff_shapeOperatorsCommute, Riemannian.DCImmersedPatch.eq_of_inner_eq_of_mem_normalSpace}
  \uses{prop:dc-ch6-3-1, def:dc-ch6-3-covderiv-B, def:dc-ch6-3-constant-curvature, def:dc-ch6-2-2}
  \leanok
  \dcref{ch6:3.3}
  We say the normal bundle of an immersion is \emph{flat} if $R^\perp=0$. Assume the
  ambient space $\overline{M}$ has constant sectional curvature. Then the Ricci
  equation can be written
  $$
  \langle R^\perp(X,Y)\eta,\zeta\rangle=-\langle[S_\eta,S_\zeta]X,Y\rangle.
  $$
  Thus $R^\perp=0$ if and only if $[S_\eta,S_\zeta]=0$ for all
  $\eta,\zeta$, that is, if and only if for all $p\in M$ there exists a basis of
  $T_pM$ which diagonalizes all of the $S_\eta$ simultaneously.
\end{remark}

\begin{definition}[the covariant derivative of the second fundamental form]
  \label{def:dc-ch6-3-covderiv-B}
  \lean{Riemannian.DCImmersedPatch.secondFundTensor, Riemannian.DCImmersedPatch.secondFundTensorCovDeriv}
  \uses{def:dc-ch6-2-B, def:dc-ch6-2-induced-connection, def:dc-ch6-3-normal-connection}
  \leanok
  The second fundamental form regarded as a tensor
  $B:\mathcal{X}(M)\times\mathcal{X}(M)\times\mathcal{X}(M)^\perp\to\mathbb{R}$
  is $B(X,Y,\eta)=\langle B(X,Y),\eta\rangle$, with covariant derivative

  $$
  (\overline{\nabla}_X B)(Y,Z,\eta)=X(B(Y,Z,\eta))-B(\nabla_X Y,Z,\eta)-B(Y,\nabla_X Z,\eta)-B(Y,Z,\nabla_X^\perp\eta).
  $$
\end{definition}

\begin{lemma}[reduction of the covariant derivative of $B$]
  \label{lem:dc-ch6-3-covderiv-reduce}
  \lean{Riemannian.DCImmersedPatch.secondFundTensorCovDeriv_eq}
  \uses{def:dc-ch6-3-covderiv-B, lem:dc-ch6-3-weingarten}
  \leanok
  For $\eta$ normal,

  $$
  (\overline{\nabla}_X B)(Y,Z,\eta)=\langle\nabla_X^\perp(B(Y,Z)),\eta\rangle-\langle B(\nabla_X Y,Z),\eta\rangle-\langle B(Y,\nabla_X Z),\eta\rangle.
  $$
\end{lemma}
\begin{proof}
  \leanok
  \sketch
  By compatibility of $\overline{\nabla}$ with the metric,
  $X\langle B(Y,Z),\eta\rangle
  =\langle\overline{\nabla}_X(B(Y,Z)),\eta\rangle
  +\langle B(Y,Z),\overline{\nabla}_X\eta\rangle$. In the first term only the
  normal part $\nabla^\perp_X(B(Y,Z))$ of $\overline{\nabla}_X(B(Y,Z))$
  survives against the normal field $\eta$; in the second only the normal part
  $\nabla^\perp_X\eta$ of $\overline{\nabla}_X\eta$ survives against the normal
  field $B(Y,Z)$. Substituting into the definition of
  $(\overline{\nabla}_X B)(Y,Z,\eta)$ cancels the term
  $\langle B(Y,Z),\nabla^\perp_X\eta\rangle$ and leaves the stated formula.
\end{proof}

\begin{proposition}[Codazzi's equation]
  \label{prop:dc-ch6-3-4}
  \lean{Riemannian.DCImmersedPatch.codazzi_equation}
  \uses{def:dc-ch6-3-covderiv-B, lem:dc-ch6-3-covderiv-reduce, lem:dc-ch6-3-decomposition, prop:dc-ch6-2-induced-connection}
  \leanok
  \dcref{ch6:3.4}
  With the notation above,

  $$
  \langle\overline{R}(X,Y)Z,\eta\rangle=(\overline{\nabla}_Y B)(X,Z,\eta)-(\overline{\nabla}_X B)(Y,Z,\eta).
  $$
\end{proposition}
\begin{proof}
  \leanok
  \sketch
  Observe first that, by \cref{lem:dc-ch6-3-covderiv-reduce},

  $$
  (\overline{\nabla}_X B)(Y,Z,\eta)=X\langle B(Y,Z),\eta\rangle-\langle B(\nabla_X Y,Z),\eta\rangle-\langle B(Y,\nabla_X Z),\eta\rangle-\langle B(Y,Z),\nabla_X^\perp\eta\rangle=\langle\nabla_X^\perp(B(Y,Z)),\eta\rangle-\langle B(\nabla_X Y,Z),\eta\rangle-\langle B(Y,\nabla_X Z),\eta\rangle.
  $$

  Consider expression (6) of \cref{lem:dc-ch6-3-decomposition} and multiply both sides by
  $\eta$. Taking the remark above into account,
  
  $$
  \langle\overline{R}(X,Y)Z,\eta\rangle=\langle B(Y,\nabla_X Z),\eta\rangle+\langle\nabla_Y^\perp B(X,Z),\eta\rangle-\langle B(X,\nabla_Y Z),\eta\rangle-\langle\nabla_X^\perp B(Y,Z),\eta\rangle+\langle B(\nabla_X Y,Z),\eta\rangle-\langle B(\nabla_Y X,Z),\eta\rangle=-(\overline{\nabla}_X B)(Y,Z,\eta)+(\overline{\nabla}_Y B)(X,Z,\eta),
  $$
  
  which is Codazzi's equation.
\end{proof}

\begin{remark}[Codazzi for constant curvature ambient and hypersurfaces]
  \label{rem:dc-ch6-3-5}
  \lean{Riemannian.DCImmersedPatch.codazzi_symm_of_isConstantCurvature, Riemannian.DCImmersedPatch.normalCov_eq_zero_of_unit_spanning, Riemannian.DCImmersedPatch.codazzi_hypersurface}
  \uses{prop:dc-ch6-3-4, def:dc-ch6-3-constant-curvature, def:dc-ch6-3-normal-connection}
  \leanok
  \dcref{ch6:3.5}
  If the ambient space $\overline{M}$ has constant sectional curvature, the Codazzi
  equation reduces to
  $$
  (\overline{\nabla}_X B)(Y,Z,\eta)=(\overline{\nabla}_Y B)(X,Z,\eta).
  $$
  If, in addition, the codimension of the immersion is $1$ and $\eta$ is chosen as a
  local unit normal field, then
  $\nabla_X^\perp\eta=0$, hence
  $$
  (\overline{\nabla}_X B)(Y,Z,\eta)=X\langle S_\eta(Y),Z\rangle-\langle S_\eta(\nabla_X Y),Z\rangle-\langle S_\eta(Y),\nabla_X Z\rangle=\langle\nabla_X(S_\eta(Y)),Z\rangle-\langle S_\eta(\nabla_X Y),Z\rangle.
  $$
  Therefore, in this case, the Codazzi equation can be written
  $$
  \nabla_X(S_\eta(Y))-\nabla_Y(S_\eta(X))=S_\eta([X,Y]).
  $$
\end{remark}

\begin{lemma}[pointwise geodesic-field criterion]
  \label{lem:dc-ch6-2-9-field}
  \source{docarmo:page-0147}
  \dcref{ch6:2.9}
  \uses{def:dc-ch6-2-8-geodesic, def:dc-ch6-2-9-geodesic-field,
    lem:dc-ch6-2-9-existence, def:dc-ch6-2-B, def:dc-ch6-2-2,
    lem:dc-ch6-2-locality, prop:dc-ch6-2-1}
  \lean{Riemannian.DCImmersedPatch.isGeodesicAt_iff_forall_isGeodesicFieldAt,
    Riemannian.DCImmersedPatch.secondFundQuadAt_eq_metricInner_cov,
    Riemannian.DCImmersedPatch.IsTotallyGeodesic.cov_apply_eq_zero}
  \leanok
  An immersion $f:M\to\overline{M}$ is geodesic at $p\in M$ if and only if every
  geodesic field $X$ at $p$
  (\cref{def:dc-ch6-2-9-geodesic-field}) satisfies
  $(\overline{\nabla}_X X)(p)=0$. The two readings agree through the displayed
  identity of the proof: for any tangent field $X$ with $X(p)=x$ and any
  $\eta\in(T_pM)^\perp$,
  $$
  H_\eta(x,x)=\langle\overline{\nabla}_X X,\eta\rangle(p).
  $$
\end{lemma}
\begin{proof}
  \leanok
  For any tangent field $X$ with $X(p)=x$ and $\eta\in(T_pM)^\perp$, the Gauss
  decomposition $\overline{\nabla}_X X=\nabla_X X+B(X,X)$ paired against $\eta$
  gives (do Carmo computes the same identity through the Weingarten pairing
  $\langle N,\overline{\nabla}_X X\rangle=-\langle\overline{\nabla}_X N,X\rangle$
  for a normal extension $N$ of $\eta$)
  $$
  H_\eta(x,x)=\langle B(x,x),\eta\rangle=\langle\overline{\nabla}_X X,\eta\rangle(p),
  $$
  since the tangential part $\nabla_X X$ pairs to zero against $\eta$.

  ($\Rightarrow$) If $f$ is geodesic at $p$ then $B$ vanishes at $p$, so for a
  geodesic field $X$ at $p$ the ambient acceleration
  $(\overline{\nabla}_X X)(p)=(\nabla_X X)(p)+B(x,x)=0+0=0$.

  ($\Leftarrow$) Let $\eta\in(T_pM)^\perp$ and $x\in T_pM$. By
  \cref{lem:dc-ch6-2-9-existence} there is a geodesic field $X$ at $p$ with
  $X(p)=x$; by hypothesis $(\overline{\nabla}_X X)(p)=0$, hence
  $H_\eta(x,x)=\langle\overline{\nabla}_X X,\eta\rangle(p)=0$. Thus every
  $\mathit{II}_\eta$ vanishes at $p$, i.e.\ $f$ is geodesic at $p$. $\square$
\end{proof}
